% ==============================================================================
% Formation C# & SQL — Module 001 : Les Bases
% Fichier : src/P1_08_a.tex
% Sujet   : Chapitre 8 : Manipulation des types de base (Chaînes et Nombres)
% ==============================================================================

\newpage
\section{Manipulation des Types de Base (Chaînes et Nombres)}
\marginpar{\tiny\color{gray}P1\_08\_a.tex}

En C\#, le type \texttt{string} (alias de la classe \texttt{System.String}) représente une séquence de caractères Unicode. Une caractéristique fondamentale à assimiler en ingénierie logicielle est l'\textbf{immutabilité} des objets \texttt{string} : une fois créée en mémoire, une chaîne de caractères ne peut plus être modifiée. Toute opération de transformation (mise en majuscules, remplacement, découpage) instancie et retourne une \textit{nouvelle} chaîne en mémoire, laissant la chaîne d'origine intacte.

\subsection{Propriétés et Méthodes de Base}

\begin{lstlisting}[style=csharpstyle]
string texte = " Connexion au serveur... ";

// Propriété Length : Obtient la longueur totale (incluant les espaces)
int longueur = texte.Length; // 25

// Trim() : Supprime les espaces blancs (espaces, tabulations, sauts de ligne) aux extrémités
string clean = texte.Trim(); // "Connexion au serveur..."

// Transformation de la casse
string majuscule = clean.ToUpper(); // "CONNEXION AU SERVEUR..."
string minuscule = clean.ToLower(); // "connexion au serveur..."
\end{lstlisting}

\subsection{Recherche et Vérification}

Pour vérifier le contenu d'une chaîne sans nécessairement l'extraire, C\# propose des méthodes d'évaluation retournant un booléen ou un index.

\begin{lstlisting}[style=csharpstyle]
string log = "[ERROR] Timeout de la base de données";

// Vérifications booléennes
bool estErreur = log.StartsWith("[ERROR]"); // true
bool concerneBDD = log.EndsWith("données"); // true
bool contientTimeout = log.Contains("Timeout"); // true

// Recherche d'index (retourne la position zéro-indexée, ou -1 si introuvable)
int position = log.IndexOf("Timeout"); // 8
\end{lstlisting}

\subsection{Découpage, Concaténation et Remplacement}

La manipulation de données brutes (comme des fichiers CSV ou des logs) requiert souvent d'extraire, d'isoler ou de remplacer des segments de texte.

\begin{lstlisting}[style=csharpstyle]
// Substring(startIndex, length) : Extrait un segment précis
string identifiant = "USER_1024_ADMIN";
string idNum = identifiant.Substring(5, 4); // Extrait "1024"

// Replace(oldValue, newValue) : Remplace toutes les occurrences
string requeteSql = "SELECT * FROM [Table]";
string requeteSecurisee = requeteSql.Replace("[Table]", "Utilisateurs");

// Split() : Divise une chaîne en tableau (ex: lecture de CSV)
string ligneCsv = "Alice,Dupont,30,Paris";
string[] colonnes = ligneCsv.Split(','); 
// colonnes = ["Alice", "Dupont", "30", "Paris"]

// Join() : Recompose un tableau en une chaîne avec un séparateur
string nouvelleLigne = string.Join(";", colonnes); 
// "Alice;Dupont;30;Paris"
\end{lstlisting}

\subsection{Insertion et Suppression}

\begin{lstlisting}[style=csharpstyle]
string chemin = "C:\\dossier\\fichier.txt";

// Insert(index, valeur) : Insère une chaîne à une position précise
string cheminBackup = chemin.Insert(18, "_backup"); 
// "C:\dossier\fichier_backup.txt"

// Remove(index, length) : Supprime une partie de la chaîne
string nomCourt = cheminBackup.Remove(0, 11); 
// "fichier_backup.txt"
\end{lstlisting}

\subsection{Interpolation et Optimisation (\texttt{StringBuilder})}

\subsubsection{Interpolation de chaînes (\texttt{\$""})}
L'interpolation, introduite en C\# 6, permet d'injecter directement des expressions ou variables dans une chaîne à l'aide d'accolades \texttt{\{\}}. C'est l'approche la plus lisible pour la construction de chaînes dynamiques, bien supérieure à la concaténation par l'opérateur \texttt{+}.

\begin{lstlisting}[style=csharpstyle]
string serveur = "192.168.1.10";
int port = 443;
// Le préfixe $ active l'interpolation
string uri = $"https://{serveur}:{port}/api/v1/status";
\end{lstlisting}

\subsubsection{Optimisation avec \texttt{StringBuilder}}
À cause de l'immutabilité du type \texttt{string}, une boucle qui concatène des milliers de chaînes via \texttt{+=} provoquera autant d'allocations mémoires superflues, saturant le \textit{Garbage Collector}. Dans les contextes de forte performance ou de modifications intensives, il est impératif d'utiliser la classe \texttt{StringBuilder} (namespace \texttt{System.Text}).

\begin{lstlisting}[style=csharpstyle]
using System.Text;

StringBuilder sb = new StringBuilder();
for (int i = 0; i < 10000; i++)
{
    // Ne recrée pas de nouvel objet à chaque tour, modifie le buffer existant
    sb.Append($"Ligne {i} | ");
}
string resultatFinal = sb.ToString();
\end{lstlisting}

\begin{tcolorbox}[colback=backcolour,colframe=keywordcolor,title=\textbf{Pièges Courants \& Erreurs}]
\begin{itemize}
    \item \textbf{Immutabilité ignorée} : Écrire \texttt{texte.Trim();} seul sans récupérer le résultat (\texttt{texte = texte.Trim();}) n'aura aucun effet sur la variable \texttt{texte}.
    \item \textbf{NullReferenceException} : Appeler \texttt{Contains} ou \texttt{ToLower} sur une variable \texttt{string} valant \texttt{null} provoquera un crash instantané. Il faut toujours vérifier la nullité avec \texttt{string.IsNullOrEmpty(texte)} ou \texttt{string.IsNullOrWhiteSpace(texte)} avant manipulation.
    \item \textbf{Concaténation dans des boucles} : Utiliser l'opérateur \texttt{+=} sur une \texttt{string} dans une boucle de plus d'une dizaine d'itérations est considéré comme un défaut d'architecture (\textit{Code Smell}) pour les performances mémoire.
\end{itemize}
\end{tcolorbox}

\subsection{Manipulation des Nombres et Mathématiques}

La classe statique \texttt{System.Math} fournit l'ensemble des opérations mathématiques fondamentales nécessaires en ingénierie logicielle. Parallèlement, les types primitifs numériques disposent de méthodes de parsing sécurisées.

\begin{lstlisting}[style=csharpstyle]
// Méthodes mathématiques de base
int absolu = Math.Abs(-5); // 5
int maximum = Math.Max(10, 20); // 20
double puissance = Math.Pow(2, 3); // 8 (2 au cube)
double arrondi = Math.Round(5.5); // 6

// Limitation d'une valeur dans un intervalle (Clamp)
int valeurRestreinte = Math.Clamp(15, 0, 10); // 10
\end{lstlisting}

\subsubsection{Conversion et Parsing}
La conversion d'une chaîne de caractères vers un entier est une opération risquée qui doit toujours être sécurisée.

\begin{lstlisting}[style=csharpstyle]
string saisie = "42";

// Parse : Lève une exception (FormatException) si la chaîne est invalide
// Toujours passer CultureInfo.InvariantCulture pour un parsing numérique déconnecté de la locale du système
int nombre = int.Parse(saisie, System.Globalization.CultureInfo.InvariantCulture);

// TryParse : Approche sécurisée recommandé en prod (ne lève pas d'exception, retourne un booléen)
if (int.TryParse("500.00", System.Globalization.NumberStyles.Any, System.Globalization.CultureInfo.InvariantCulture, out int resultatSecurise))
{
    Console.WriteLine($"Conversion réussie : {resultatSecurise}");
}
else
{
    Console.WriteLine("Échec de la conversion (format ou culture non conforme).");
}
\end{lstlisting}

\subsection{Synthèse}
\begin{itemize}
    \item \textbf{Propriétés fondamentales} : Les objets \texttt{string} sont immuables. Chaque modification génère une nouvelle allocation mémoire.
    \item \textbf{Évaluation} : Utiliser \texttt{StartsWith}, \texttt{Contains} et \texttt{IndexOf} pour analyser le contenu sans créer de nouvelles chaînes.
    \item \textbf{Formatage dynamique} : Privilégier systématiquement l'interpolation (\texttt{\$""}) pour la lisibilité du code.
    \item \textbf{Performance} : Utiliser exclusivement \texttt{StringBuilder} lors de la construction itérative de chaînes dans des boucles.
\end{itemize}
